\documentclass[10pt,a4paper]{amsart}
\usepackage{amsmath,amssymb,amsthm,graphicx}
\usepackage{verbatim}
\newenvironment{reference}{\comment}{\endcomment}
%\newenvironment{reference}{}{}
\newenvironment{slogan}{\comment}{\endcomment}
\newenvironment{history}{\comment}{\endcomment}

% For commutative diagrams we use Xy-pic
\usepackage[all]{xy}

% We use 2cell for 2-commutative diagrams.
\xyoption{2cell}
% Theorem environments.
%
\theoremstyle{plain}
\newtheorem{theorem}[subsection]{Theorem}
\newtheorem{proposition}[subsection]{Proposition}
\newtheorem{lemma}[subsection]{Lemma}

\theoremstyle{definition}
\newtheorem{definition}[subsection]{Definition}
\newtheorem{example}[subsection]{Example}
\newtheorem{exercise}[subsection]{Exercise}
\newtheorem{situation}[subsection]{Situation}

\theoremstyle{remark}
\newtheorem{remark}[subsection]{Remark}
\newtheorem{remarks}[subsection]{Remarks}

\numberwithin{equation}{subsection}

% Macros
%
\def\lim{\mathop{\mathrm{lim}}\nolimits}
\def\colim{\mathop{\mathrm{colim}}\nolimits}
\def\Spec{\mathop{\mathrm{Spec}}}
\def\Hom{\mathop{\mathrm{Hom}}\nolimits}
\def\Ext{\mathop{\mathrm{Ext}}\nolimits}
\def\SheafHom{\mathop{\mathcal{H}\!\mathit{om}}\nolimits}
\def\SheafExt{\mathop{\mathcal{E}\!\mathit{xt}}\nolimits}
\def\Sch{\mathit{Sch}}
\def\Mor{\mathop{\mathrm{Mor}}\nolimits}
\def\Ob{\mathop{\mathrm{Ob}}\nolimits}
\def\Sh{\mathop{\mathit{Sh}}\nolimits}
\def\NL{\mathop{N\!L}\nolimits}
\def\CH{\mathop{\mathrm{CH}}\nolimits}
\def\proetale{{pro\text{-}\acute{e}tale}}
\def\etale{{\acute{e}tale}}
\def\QCoh{\mathit{QCoh}}
\def\Ker{\mathop{\mathrm{Ker}}}
\def\Im{\mathop{\mathrm{Im}}}
\def\Coker{\mathop{\mathrm{Coker}}}
\def\Coim{\mathop{\mathrm{Coim}}}

% Boxtimes
%
\DeclareMathSymbol{\boxtimes}{\mathbin}{AMSa}{"02}

%
% Macros for moduli stacks/spaces
%
\def\QCohstack{\mathcal{QC}\!\mathit{oh}}
\def\Cohstack{\mathcal{C}\!\mathit{oh}}
\def\Spacesstack{\mathcal{S}\!\mathit{paces}}
\def\Quotfunctor{\mathrm{Quot}}
\def\Hilbfunctor{\mathrm{Hilb}}
\def\Curvesstack{\mathcal{C}\!\mathit{urves}}
\def\Polarizedstack{\mathcal{P}\!\mathit{olarized}}
\def\Complexesstack{\mathcal{C}\!\mathit{omplexes}}
% \Pic is the operator that assigns to X its picard group, usage \Pic(X)
% \Picardstack_{X/B} denotes the Picard stack of X over B
% \Picardfunctor_{X/B} denotes the Picard functor of X over B
\def\Pic{\mathop{\mathrm{Pic}}\nolimits}
\def\Picardstack{\mathcal{P}\!\mathit{ic}}
\def\Picardfunctor{\mathrm{Pic}}
\def\Deformationcategory{\mathcal{D}\!\mathit{ef}}
\usepackage{cleveref}
\crefname{originallemma}{original lemma}{original lemmas}
\Crefname{originallemma}{Original lemma}{Original lemmas}
\makeatletter
\newlabel{chow-lemma-prepare-flat-pullback-rational-equivalence}{{0EPH}{source}{}{sourcepointer-1}{}}
\newlabel{chow-lemma-prepare-flat-pullback-rational-equivalence@cref}{{[originallemma][0][]0EPH}{[1][0][]source}{}{}{}}
\makeatother
\crefname{originalreference}{original reference}{original references}
\Crefname{originalreference}{Original reference}{Original references}
\makeatletter
\newlabel{lemma-compose-flat-pullback}{{02RD}{source}{}{sourcepointer-2}{}}
\newlabel{lemma-compose-flat-pullback@cref}{{[originalreference][0][]02RD}{[1][0][]source}{}{}{}}
\makeatother
\crefname{originalreference}{original reference}{original references}
\Crefname{originalreference}{Original reference}{Original references}
\makeatletter
\newlabel{lemma-etale-pullback}{{0EPY}{source}{}{sourcepointer-3}{}}
\newlabel{lemma-etale-pullback@cref}{{[originalreference][0][]0EPY}{[1][0][]source}{}{}{}}
\makeatother
\crefname{originalreference}{original reference}{original references}
\Crefname{originalreference}{Original reference}{Original references}
\makeatletter
\newlabel{lemma-etale-pullback-principal-divisor}{{0EQ2}{source}{}{sourcepointer-4}{}}
\newlabel{lemma-etale-pullback-principal-divisor@cref}{{[originalreference][0][]0EQ2}{[1][0][]source}{}{}{}}
\makeatother
\crefname{originalreference}{original reference}{original references}
\Crefname{originalreference}{Original reference}{Original references}
\makeatletter
\newlabel{lemma-flat-inverse-image-dimension}{{02R8}{source}{}{sourcepointer-5}{}}
\newlabel{lemma-flat-inverse-image-dimension@cref}{{[originalreference][0][]02R8}{[1][0][]source}{}{}{}}
\makeatother
\crefname{originalreference}{original reference}{original references}
\Crefname{originalreference}{Original reference}{Original references}
\makeatletter
\newlabel{lemma-flat-pullback-proper-pushforward}{{02RG}{source}{}{sourcepointer-6}{}}
\newlabel{lemma-flat-pullback-proper-pushforward@cref}{{[originalreference][0][]02RG}{[1][0][]source}{}{}{}}
\makeatother
\crefname{originalreference}{original reference}{original references}
\Crefname{originalreference}{Original reference}{Original references}
\makeatletter
\newlabel{lemma-proper-image}{{0EEG}{source}{}{sourcepointer-7}{}}
\newlabel{lemma-proper-image@cref}{{[originalreference][0][]0EEG}{[1][0][]source}{}{}{}}
\makeatother
\crefname{originalreference}{original reference}{original references}
\Crefname{originalreference}{Original reference}{Original references}
\makeatletter
\newlabel{situation-setup}{{02QL}{source}{}{sourcepointer-8}{}}
\newlabel{situation-setup@cref}{{[originalreference][0][]02QL}{[1][0][]source}{}{}{}}
\makeatother
\crefname{originalsection}{original section}{original sections}
\Crefname{originalsection}{Original section}{Original sections}
\makeatletter
\newlabel{spaces-morphisms-section-rational-maps}{{0EML}{source}{}{sourcepointer-9}{}}
\newlabel{spaces-morphisms-section-rational-maps@cref}{{[originalsection][0][]0EML}{[1][0][]source}{}{}{}}
\makeatother
\title{Flat pullback of a principal divisor equals the multiplicity weighted sum over the irreducible components of the source}
\author{The Stacks Project Authors}
\date{stacks-project:spaces-chow.tex:local-snapshot-20261012; GNU FDL 1.2 or later}
\begin{document}
\maketitle
\noindent\textit{Editorial note.} All delivered mathematical text is the author's own, copied from the pinned Stacks Project chapter file spaces-chow.tex; no statement or proof step was rewritten, reordered or removed.\par\medskip
\noindent\textit{Editorial note.} The counted proof is the author's complete argument: a two-stage etale descent to the scheme case, with the pullback of the divisor computed in two ways and the multiplicities matched through the equality of flat pullbacks of the fundamental class\par\medskip
\noindent\textit{Editorial note.} The label lemma-prepare-flat-pullback-rational-equivalence names the counted statement, and the proof environment immediately following it in the source is delivered verbatim as its proof.\par\medskip
\noindent\textit{Editorial note.} The counted statement is carried by the author's own lemma environment in the chapter source.\par\medskip
\noindent\textit{Allowed standard prerequisites (editorial).}\begin{itemize}
\item The etale pullback formula for principal divisors, and the principle that a cycle equality may be checked after pulling back along a surjective etale morphism
\item Flat pullback and proper pushforward of cycles, the compatibility of flat pullback with proper pushforward (projection formula) and the composition of flat pullbacks
\item The scheme case of the statement, and the multiplicity of an irreducible component in a cycle together with the scheme theoretic image of a component
\end{itemize}\medskip
\ifdefined\hypertarget\hypertarget{sourcepointer-1}{}\fi
\noindent\textit{Source reference.} Original lemma 0EPH: sources/stacks-chow/tags.txt, source lines 18377--18377. Prerequisite pointer invoked at proof line(s) 2082; it names a standard fact of a sibling chapter that the counted argument uses, and every step of the counted proof built on it is retained verbatim in the delivered body.\par\medskip
\ifdefined\hypertarget\hypertarget{sourcepointer-2}{}\fi
\noindent\textit{Source reference.} Original reference 02RD: sources/stacks-spaces-chow/pinned-original-source.tex, source lines 1347--1347. Prerequisite pointer invoked at proof line(s) 2103; it names a standard fact of the same chapter that the counted argument uses, and every step of the counted proof built on it is retained verbatim in the delivered body.\par\medskip
\ifdefined\hypertarget\hypertarget{sourcepointer-3}{}\fi
\noindent\textit{Source reference.} Original reference 0EPY: sources/stacks-spaces-chow/pinned-original-source.tex, source lines 1332--1332. Prerequisite pointer invoked at proof line(s) 2067; it names a standard fact of the same chapter that the counted argument uses, and every step of the counted proof built on it is retained verbatim in the delivered body.\par\medskip
\ifdefined\hypertarget\hypertarget{sourcepointer-4}{}\fi
\noindent\textit{Source reference.} Original reference 0EQ2: sources/stacks-spaces-chow/pinned-original-source.tex, source lines 1680--1680. Prerequisite pointer invoked at proof line(s) 2069; it names a standard fact of the same chapter that the counted argument uses, and every step of the counted proof built on it is retained verbatim in the delivered body.\par\medskip
\ifdefined\hypertarget\hypertarget{sourcepointer-5}{}\fi
\noindent\textit{Source reference.} Original reference 02R8: sources/stacks-spaces-chow/pinned-original-source.tex, source lines 1149--1149. Prerequisite pointer invoked at proof line(s) 2054; it names a standard fact of the same chapter that the counted argument uses, and every step of the counted proof built on it is retained verbatim in the delivered body.\par\medskip
\ifdefined\hypertarget\hypertarget{sourcepointer-6}{}\fi
\noindent\textit{Source reference.} Original reference 02RG: sources/stacks-spaces-chow/pinned-original-source.tex, source lines 1498--1498. Prerequisite pointer invoked at proof line(s) 2074; it names a standard fact of the same chapter that the counted argument uses, and every step of the counted proof built on it is retained verbatim in the delivered body.\par\medskip
\ifdefined\hypertarget\hypertarget{sourcepointer-7}{}\fi
\noindent\textit{Source reference.} Original reference 0EEG: sources/stacks-spaces-chow/pinned-original-source.tex, source lines 820--820. Prerequisite pointer invoked at proof line(s) 2097; it names a standard fact of the same chapter that the counted argument uses, and every step of the counted proof built on it is retained verbatim in the delivered body.\par\medskip
\ifdefined\hypertarget\hypertarget{sourcepointer-8}{}\fi
\noindent\textit{Source reference.} Original reference 02QL: sources/stacks-spaces-chow/pinned-original-source.tex, source lines 264--264. Prerequisite pointer invoked at proof line(s) 2040; it names a standard fact of the same chapter that the counted argument uses, and every step of the counted proof built on it is retained verbatim in the delivered body.\par\medskip
\ifdefined\hypertarget\hypertarget{sourcepointer-9}{}\fi
\noindent\textit{Source reference.} Original section 0EML: sources/stacks-chow/tags.txt, source lines 18311--18311. Prerequisite pointer invoked at proof line(s) 2056; it names a standard fact of a sibling chapter that the counted argument uses, and every step of the counted proof built on it is retained verbatim in the delivered body.\par\medskip
\begin{lemma}
\label{lemma-prepare-flat-pullback-rational-equivalence}
In Situation \ref{situation-setup} let $X, Y/B$ be good.
Assume $Y$ integral with $\dim_\delta(Y) = k$.
Let $f : X \to Y$ be a flat morphism of
relative dimension $r$. Then for $g \in R(Y)^*$ we have
$$
f^*\text{div}_Y(g) =
\sum m_{X', X} (X' \to X)_*\text{div}_{X'}(g \circ f|_{X'})
$$
as $(k + r - 1)$-cycles on $X$ where the sum is over the irreducible
components $X'$ of $X$ and $m_{X', X}$ is the multiplicity of $X'$ in $X$.
\end{lemma}

\begin{proof}
Observe that any irreducible component of $X$ dominates $Y$
(Lemma \ref{lemma-flat-inverse-image-dimension})
and hence the composition $g \circ f|_{X'}$ is defined
(Morphisms of Spaces, Section \ref{spaces-morphisms-section-rational-maps}).
We will reduce this to the case of schemes. Choose a scheme $V$ and a
surjective \'etale morphism $V \to Y$. Choose a scheme $U$ and a surjective
\'etale morphism $U \to V \times_Y X$. Picture
$$
\xymatrix{
U \ar[r]_a \ar[d]_h & X \ar[d]^f \\
V \ar[r]^b & Y
}
$$
Since $a$ is surjective and \'etale it follows from
Lemma \ref{lemma-etale-pullback} that it suffices to prove the
equality of cycles after pulling back by $a$.
We can use Lemma \ref{lemma-etale-pullback-principal-divisor} to write
$$
b^*\text{div}_Y(g) = \sum (V' \to V)_*\text{div}_{V'}(g \circ b|_{V'})
$$
where the sum is over the irreducible components $V'$ of $V$.
Using Lemma \ref{lemma-flat-pullback-proper-pushforward} we find
$$
h^*b^*\text{div}_Y(g) =
\sum (V' \times_V U \to U)_*(h')^*\text{div}_{V'}(g \circ b|_{V'})
$$
where $h' : V' \times_V U \to V'$ is the projection.
We may apply the lemma in the case of schemes
(Chow Homology, Lemma
\ref{chow-lemma-prepare-flat-pullback-rational-equivalence})
to the morphism $h' : V' \times_V U \to V'$ to see that we have
$$
(h')^*\text{div}_{V'}(g \circ b|_{V'}) =
\sum
m_{U', V' \times_V U}
(U' \to V' \times_V U)_*\text{div}_{U'}(g \circ b|_{V'} \circ h'|_{U'})
$$
where the sum is over the irreducible components $U'$ of $V' \times_V U$.
Each $U'$ occurring in this sum is an irreducible component of $U$
and conversely every irreducible component $U'$ of $U$ is an
irreducible component of $V' \times_V U$
for a unique irreducible component $V' \subset V$.
Given an irreducible component $U' \subset U$,
denote $\overline{a(U')} \subset X$ the ``image'' in $X$
(Lemma \ref{lemma-proper-image}); this
is an irreducible component of $X$ for example by
Lemma \ref{lemma-flat-inverse-image-dimension}.
The muplticity $m_{U', V' \times_V U}$ is equal to
the multiplicity $m_{\overline{a(U')}, X}$.
This follows from the equality $h^*a^*[Y] = b^*f^*[Y]$
(Lemma \ref{lemma-compose-flat-pullback}), the definitions, and
Lemma \ref{lemma-etale-pullback}.
Combining all of what we just said we obtain
$$
a^*f^*\text{div}_Y(g) =
h^*b^*\text{div}_Y(g) =
\sum m_{\overline{a(U')}, X}
(U' \to U)_*\text{div}_{U'}(g \circ (f \circ a)|_{U'})
$$
Next, we analyze what happens with the right hand side of the
formula in the statement of the lemma if we pullback by $a$.
First, we use Lemma \ref{lemma-flat-pullback-proper-pushforward}
to get
$$
a^*\sum m_{X', X} (X' \to X)_*\text{div}_{X'}(g \circ f|_{X'}) =
\sum m_{X', X} (X' \times_X U \to U)_*(a')^*\text{div}_{X'}(g \circ f|_{X'})
$$
where $a' : X' \times_X U \to X'$ is the projection.
By Lemma \ref{lemma-etale-pullback-principal-divisor}
we get
$$
(a')^*\text{div}_{X'}(g \circ f|_{X'}) =
\sum (U' \to X' \times_X U)_*\text{div}_{U'}(g \circ (f \circ a)|_{U'})
$$
where the sum is over the irreducible components $U'$ of $X' \times_X U$.
These $U'$ are irreducible components of $U$ and in fact are exactly the
irreducible components of $U$ such that $\overline{a(U')} = X'$.
Comparing with what we obtained above we conclude.
\end{proof}

\end{document}
